
\begin{proposition}[The genus-two node alternative is excluded]\label{cdp:st:C.2.CR.3}
In the genus-two situation of \ref{cdp:st:C.2.CR.2}, the component
of $\Gamma$ through $q=h(C)$ cannot have generic
transversal an ordinary node. Consequently the
remaining genus-two alternative has an ordinary cusp
at its generic point, and exactly one strict conductor
curve of coefficient two and degree one lies above
that component.
\end{proposition}
\begin{proof}
We first record an additional consequence of the
specified boundary blowups in C.2.CR.1: every
coefficient-one component of $D_{\rm red}$ is a leaf.
This holds on the initial section-and-fibres star.
A smooth boundary blowup creates coefficient one
only when the old coefficient is two, and creates
it as a leaf. A node blowup could create coefficient
one only from an edge $(1,1)$, which (C.CR.2)
excludes. Subsequent blowups preserve the valence
one of an existing coefficient-one vertex: at its
unique boundary node they merely subdivide that edge,
and at a smooth point its new exceptional coefficient
is zero.

Let $E_0$ be the connected exceptional fibre containing
$C$. In the boundary tree, the first nonexceptional
vertex on every path from $C$ meets $E_0$, so its image
passes through $p$. Its further image is therefore
a component of $\Gamma$ through $q$. By \ref{cdp:st:C.2.CR.2}
there is only one such component, denoted $G$.

Suppose its generic transversal were a node.
Every strict conductor curve above $G$ then has
coefficient one. Thus each first nonexceptional
vertex just found is a leaf. Every strict conductor
curve is consequently such a leaf, attached at a
single point to $E_0$. No additional exceptional
tail can meet it, since that would raise its
boundary valence. The reduced conductor on $T$
is therefore a union of rational curves whose only
possible multibranch point is $p$; every component
passes through $p$, is smooth away from $p$, and
has only one branch at $p$.

All their images pass through $q$, where $\Gamma$
is smooth, so the conductor image has just the one
global component $G$. It also has no multibranch
point elsewhere. Indeed, at $x\ne q$ each point
of the normalization over $x$ receives just one
local branch of $G$. If $G$ had two distinct local
branches there, their receiving normalization
points would be distinct. The corresponding
irreducible hypersurface branches of $S$ must
intersect in a curve, because their distinct
local equations generate a height-two ideal in
the regular local ring of the threefold. That
intersection curve is a conductor-image branch
received by both normalization points, contradicting
their different singleton branch sets.

The normalization of $G$ is rational and all its
germs are unibranch. The normalization map is a
continuous bijection from a compact space to a
Hausdorff space and hence a homeomorphism.
Thus $\pi_1(G)=1$, contradicting (C.CR.1).
The generic cusp alternative is the only one left.
Its transverse normalization is unibranch with
conductor exponent two, so it contributes exactly
one strict coefficient-two curve, of degree one
over $G$.
\end{proof}


